% =============================================================================
% 11-parallel-pipeline-hazards-notes.tex -- Lecture Notes for Week 11
% DERIVED ARTIFACT. Source of truth:
% Slides/CS401/11-parallel-pipeline-hazards.tex.
% The Notes reorganize the deck by topic while preserving its content,
% notation, worked examples, and citations.
%
% Anchor reading:
%   * Hamacher, Computer Organization and Embedded Systems, Chs. 5 and 10.
%     These are chapter-level course anchors; the current repository index does
%     not page-verify these chapters.
%   * Stallings, Computer Organization and Architecture, 11th ed., S16.4
%     (Instruction Pipelining, PDF pp. 667--685) and S20.1 (PDF p. 847).
%     The repository keeps the historical Stallings2015 bibliography key.
%
% Compile from Notes/CS401 on Windows/MiKTeX:
%   TEXINPUTS="../../Preambles;" xelatex -interaction=nonstopmode 11-parallel-pipeline-hazards-notes.tex
%   BIBINPUTS="../..;" bibtex 11-parallel-pipeline-hazards-notes
%   TEXINPUTS="../../Preambles;" xelatex -interaction=nonstopmode 11-parallel-pipeline-hazards-notes.tex
%   TEXINPUTS="../../Preambles;" xelatex -interaction=nonstopmode 11-parallel-pipeline-hazards-notes.tex
% =============================================================================

\documentclass[11pt]{article}
\usepackage[margin=1in]{geometry}
\usepackage{array}
\usepackage{caption}
% =============================================================================
% Preambles/header.tex — shared LaTeX/Beamer preamble for this project
%
% Use from a Beamer lecture by prepending:
%     \input{header}
% and compiling with TEXINPUTS=../Preambles:$TEXINPUTS (the /compile-latex
% skill does this automatically).
%
% PALETTE CONTRACT. Color names defined here MUST match the names in
% Quarto/theme-template.scss so Beamer and Quarto renderings use the same
% palette. The scripts/check-palette-sync.sh script enforces this.
%
% To customize: edit the HEX values below AND the matching $variable in
% Quarto/theme-template.scss, then run scripts/check-palette-sync.sh.
% =============================================================================

% -----------------------------------------------------------------------------
% Engine detection + encoding
%
% Our /compile-latex skill uses XeLaTeX, where inputenc is unnecessary and
% can produce encoding warnings. Only load inputenc under pdfTeX.
% -----------------------------------------------------------------------------
\usepackage{iftex}
\ifPDFTeX
  \usepackage[utf8]{inputenc}
\fi

\usepackage{amsmath, amssymb, amsthm}
\usepackage{booktabs}
\usepackage{graphicx}

% -----------------------------------------------------------------------------
% PALETTE — mirrors Quarto/theme-template.scss variable names.
% Load xcolor and define colors BEFORE anything that references them
% (notably hyperref, hypersetup, and the Beamer color assignments).
% -----------------------------------------------------------------------------
\usepackage{xcolor}

\definecolor{primary-blue}{HTML}{012169}      % headings, accents
\definecolor{primary-gold}{HTML}{B9975B}      % emphasis, borders
\definecolor{highlight-yellow}{HTML}{F2A900}  % markers, alerts
\definecolor{light-bg}{HTML}{E8EDF5}          % title / section backgrounds
\definecolor{jet}{HTML}{1A1A1A}               % body text

% Semantic colors (used in diagrams and annotations)
\definecolor{positive}{HTML}{15803D}          % good / observed / identified
\definecolor{negative}{HTML}{B91C1C}          % bad / problematic / bias
\definecolor{neutral}{HTML}{525252}           % reference / context

% Highlight accents (match .hi-* classes in the SCSS)
\definecolor{hi-slate}{HTML}{314F4F}
\definecolor{hi-green}{HTML}{2E8B57}
\definecolor{hi-red}{HTML}{C41E3A}

% Semantic aliases so skill templates and snippets can write
% \textcolor{accent}{...} without hard-coding "primary-blue".
\colorlet{accent}{primary-blue}
\colorlet{accent2}{primary-gold}

% -----------------------------------------------------------------------------
% Hyperref — loaded AFTER xcolor + palette so link colors resolve.
% -----------------------------------------------------------------------------
\usepackage{hyperref}
\hypersetup{colorlinks=true, linkcolor=primary-blue, urlcolor=primary-blue,
            citecolor=primary-blue, pdfborder={0 0 0}}

% -----------------------------------------------------------------------------
% Course code — set once per deck (after \input{header}, before \begin{document})
% via \coursecode{CS401}. Shown in the footer on every slide so a deck's course
% is identifiable without opening the title page. Empty by default.
% -----------------------------------------------------------------------------
\makeatletter
\newcommand{\coursecode}[1]{\gdef\@coursecodetext{#1}}
\gdef\@coursecodetext{}
\makeatother

% -----------------------------------------------------------------------------
% Beamer theme assignments (only apply under Beamer).
%
% We detect Beamer via \@ifclassloaded{beamer}, which is the documented,
% non-internal way. This must be wrapped in \makeatletter/\makeatother
% because \@ifclassloaded contains an @-catcode character.
% -----------------------------------------------------------------------------
\makeatletter
\@ifclassloaded{beamer}{%
  \usefonttheme{professionalfonts}
  \setbeamercolor{structure}{fg=primary-blue}
  \setbeamercolor{frametitle}{fg=primary-blue}
  \setbeamercolor{title}{fg=primary-blue}
  \setbeamercolor{subtitle}{fg=primary-gold}
  \setbeamercolor{author}{fg=primary-blue}
  \setbeamercolor{institute}{fg=neutral}
  \setbeamercolor{date}{fg=primary-gold}
  \setbeamercolor{itemize item}{fg=primary-blue}
  \setbeamercolor{itemize subitem}{fg=primary-gold}
  \setbeamercolor{alerted text}{fg=primary-gold}
  \setbeamercolor{block title}{fg=white, bg=primary-blue}
  \setbeamercolor{block body}{bg=light-bg}
  % Semantic block variants: block=definition/concept (blue), exampleblock=
  % worked example (green/positive), alertblock=key takeaway or caution (gold).
  % Keep to <=2 colored boxes per slide across ALL three variants (INV-7).
  \setbeamercolor{block title example}{fg=white, bg=positive}
  \setbeamercolor{block body example}{bg=light-bg}
  \setbeamercolor{block title alerted}{fg=white, bg=primary-gold}
  \setbeamercolor{block body alerted}{bg=light-bg}

  % Minimal footer: course code (left) + page X / N (right), no navigation clutter
  \setbeamertemplate{navigation symbols}{}
  \setbeamertemplate{footline}{%
    \color{neutral}\footnotesize\hspace{0.7em}\@coursecodetext\hfill
    \insertframenumber/\inserttotalframenumber\hspace{0.7em}\vspace{0.3em}}
}{}
\makeatother

% -----------------------------------------------------------------------------
% TikZ setup — libraries the snippet gallery uses
% -----------------------------------------------------------------------------
\usepackage{tikz}
\usetikzlibrary{arrows.meta, positioning, calc, decorations.pathreplacing,
                fit, shapes.geometric, backgrounds}

% Shared TikZ styles — snippets and hand-written diagrams can pick these up
% via `\begin{tikzpicture}[every node/.style={...}]` or by naming specific
% styles (e.g., `\node[dag-node] (X) at (0,0) {X};`).
\tikzset{
  % Node styles for causal DAGs and similar
  dag-node/.style = {
    draw=primary-blue, fill=light-bg, rounded corners=2pt,
    minimum width=1.2cm, minimum height=0.9cm, align=center, font=\small
  },
  decision-node/.style = {
    draw=primary-gold, fill=white, diamond, aspect=1.8,
    minimum width=1.8cm, minimum height=1.2cm, align=center, font=\small,
    inner sep=1pt
  },
  % Edge styles — observed vs counterfactual
  observed-edge/.style = {-{Latex[length=2.5mm]}, thick, draw=jet},
  counterfactual-edge/.style = {-{Latex[length=2.5mm]}, thick, draw=neutral, dashed},
  confound-edge/.style = {-{Latex[length=2.5mm]}, thick, draw=negative, dashed},
  % Data-point styles for plots
  observed-dot/.style = {fill=primary-blue, draw=primary-blue, circle, inner sep=1.2pt},
  counterfactual-dot/.style = {fill=white, draw=neutral, circle, inner sep=1.2pt}
}

% -----------------------------------------------------------------------------
% Convenience macros
% -----------------------------------------------------------------------------
\newcommand{\muted}[1]{\textcolor{neutral}{#1}}
\newcommand{\key}[1]{\textcolor{primary-gold}{\textbf{#1}}}
\newcommand{\good}[1]{\textcolor{positive}{#1}}
\newcommand{\bad}[1]{\textcolor{negative}{#1}}

% Standout frame macro: full-bleed dark background for section transitions.
% Beamer-only (uses \begin{frame}/\setbeamercolor/\usebeamerfont) -- do not
% call from an article-class file (e.g. Notes/<CODE>/*-notes.tex). \muted,
% \key, \good, \bad, and \coursecode above are plain xcolor/gdef and are
% safe to use from either class; verified when Notes/article-class support
% was added.
% Expand: \transitionslide{Your transition title}
\newcommand{\transitionslide}[1]{%
  \begingroup
  \setbeamercolor{background canvas}{bg=primary-blue}
  \begin{frame}[plain]
    \centering
    \vfill
    {\usebeamerfont{title}\color{white}\Large #1\par}
    \vfill
  \end{frame}
  \endgroup
}

% Section divider frame (Beamer-only), an alternative to \transitionslide with a
% darker two-line style: near-black (jet) background, a small white label line
% above a larger gold title, and a thin gold rule along the bottom edge.
% Expand: \sectiondivider{Section 2}{Pointers}
\newcommand{\sectiondivider}[2]{%
  \begingroup
  \setbeamercolor{background canvas}{bg=jet}
  \begin{frame}[plain]
    \vspace*{3.0cm}
    {\color{white}\ttfamily\large #1\par}
    \vspace{0.5cm}
    {\color{primary-gold}\LARGE #2\par}
    \begin{tikzpicture}[remember picture, overlay]
      \draw[primary-gold, line width=2.5pt]
        ($(current page.south west)+(0,0.1)$) -- ($(current page.south east)+(0,0.1)$);
    \end{tikzpicture}
  \end{frame}
  \endgroup
}

% -----------------------------------------------------------------------------
% Channel watermark — "Concepts That Click" (YouTube).
%
% Article-class files (CompetitiveExam Q/A sets, Notes, Assignments) get a
% light diagonal draft watermark on every page plus a centered footer naming
% the channel. Beamer decks get a subtle TikZ background watermark instead
% (the footline already carries the course code). Centralized here so every
% MCQ PDF is branded without per-file edits.
% -----------------------------------------------------------------------------
\makeatletter
\@ifclassloaded{article}{%
  \usepackage{draftwatermark}
  \SetWatermarkText{Concepts That Click}
  \SetWatermarkScale{1}
  \SetWatermarkFontSize{2cm}
  \SetWatermarkLightness{0.86}
  \SetWatermarkAngle{45}
  \usepackage{fancyhdr}
  \pagestyle{fancy}
  \fancyhf{}
  \fancyfoot[C]{\footnotesize\textcolor{neutral}{Concepts That Click~|~YouTube: \href{https://www.youtube.com/@concepts.thatclick}{@concepts.thatclick}}}
  \renewcommand{\headrulewidth}{0pt}
  \renewcommand{\footrulewidth}{0pt}
  \fancypagestyle{plain}{%
    \fancyhf{}%
    \fancyfoot[C]{\footnotesize\textcolor{neutral}{Concepts That Click~|~YouTube: \href{https://www.youtube.com/@concepts.thatclick}{@concepts.thatclick}}}%
    \renewcommand{\headrulewidth}{0pt}%
    \renewcommand{\footrulewidth}{0pt}%
  }
}{}
\@ifclassloaded{beamer}{%
  \setbeamertemplate{background}{%
    \begin{tikzpicture}[remember picture, overlay]
      \node[rotate=45, opacity=0.055, align=center]
        at (current page.center)
        {\Huge\textcolor{neutral}{Concepts That Click}};
    \end{tikzpicture}%
  }
}{}
\makeatother 
\coursecode{CS401}

\renewcommand{\thesection}{11.\arabic{section}}
\renewcommand{\thefigure}{11.\arabic{figure}}
\newtheorem{example}{Example}[section]

\newenvironment{definitionbox}[1]{%
  \par\noindent\textbf{Definition (#1).}\ \itshape
}{\par\vspace{0.3em}}

\title{Parallel Processing and Pipeline Hazards\\\large Lecture Notes --- Week 11}
\author{PCC CS-401: Computer Organization and Architecture}
\date{\today}

\begin{document}
\maketitle

\section{From Latency Hiding to Instruction-Level Parallelism}
\label{sec:motivation}

Week 9 used locality to hide DRAM latency behind cache, and Week 10 used
pages, frames, and storage to give a process a larger, protected address space
than physical memory alone could provide. Both mechanisms ask the same
question: how can the machine keep useful work moving while one operation is
slow? This week asks the question at the processor itself. If several useful
operations are ready at once, the processor can overlap them, but only when
the required resources and data dependencies permit the overlap.

The running instruction stream for this chapter is
\begin{center}
\texttt{LW R1,0(R2)}; \texttt{ADD R3,R1,R4}; \texttt{SUB R5,R3,R6};
\texttt{BEQ R5,R0,L}; \texttt{OR R7,R8,R9}.
\end{center}
It will let us identify dependencies and then ask which mechanism pays for
each one. The chapter proceeds in four connected topics: classifying the kind
of parallel work, constructing a five-stage pipeline, diagnosing hazards, and
choosing a resolution whose throughput gain justifies its cost.

\section{Flynn's Classification and Scales of Parallelism}
\label{sec:flynn}

\subsection{Instruction and Data Streams}

The first question is not ``how many ALUs are present?'' It is ``how many
instruction streams and data streams are moving?'' An instruction stream is
the sequence of operations being issued. A data stream is the sequence of
operands or data elements consumed by those operations. Flynn's taxonomy
classifies a computer by the number of concurrent streams; it is a taxonomy of
organization, not a ranking from bad to good. This is the standard course
treatment in Hamacher, Ch.~5 \cite{Hamacher2002_computer_organization}, while
Stallings gives a brief parallel-organization discussion in S20.1
\cite{Stallings2015_computer_organization}.

\begin{definitionbox}{Flynn's taxonomy}
Flynn's taxonomy classifies an organization by its number of instruction
streams and data streams. The four labels are SISD (single instruction,
single data), SIMD (single instruction, multiple data), MISD (multiple
instruction, single data), and MIMD (multiple instruction, multiple data).
\end{definitionbox}

The four classes can be read directly from the following table.
\begin{center}
\begin{tabular}{@{}>{\raggedright\arraybackslash}p{1.5cm}>{\raggedright\arraybackslash}p{3.0cm}>{\raggedright\arraybackslash}p{3.0cm}>{\raggedright\arraybackslash}p{5.0cm}@{}}
\toprule
\textbf{Class} & \textbf{Instruction streams} & \textbf{Data streams} &
\textbf{Typical interpretation} \\
\midrule
SISD & One & One & Conventional sequential core; one instruction acts on one data item \\
SIMD & One & Many & One operation applied across a vector or image row \\
MISD & Many & One & Rare and specialized; multiple operations inspect one stream \\
MIMD & Many & Many & Multicore, multiprocessor, or cluster with independent work \\
\bottomrule
\end{tabular}
\end{center}

One physical chip can support more than one mode at different levels. For
example, multiple cores may each run an SISD-like stream while a vector unit
executes SIMD operations. Therefore the classification should be tied to the
level being described, not assigned once to the whole chip without saying
which organization is under discussion.

\begin{example}[Classifying three workloads]
\label{ex:flynn}
Consider the following workloads.
\begin{enumerate}
  \item Add two 1024-element vectors using one vector instruction per chunk.
  \item Run four independent matrix tiles on four cores.
  \item Apply several independent filters to the same sensor stream.
\end{enumerate}
The first workload is SIMD: one operation pattern is applied to many data
elements. The second is MIMD: the cores execute independent instruction
streams over independent data regions. The third resembles MISD under a
strict stream definition because several operations inspect one input stream,
but in practice it is often implemented as MIMD workers sharing an input
buffer. The taxonomy describes the organization; the workload determines
whether that organization is useful.
\end{example}

\subsection{DLP, TLP, and ILP}

Flynn's classes describe streams, while data-level, thread-level, and
instruction-level parallelism describe where useful concurrency comes from.
\begin{definitionbox}{Data-level parallelism (DLP)}
Data-level parallelism applies the same operation to many elements, as in
vector addition or image pixels.
\end{definitionbox}
\begin{definitionbox}{Thread-level parallelism (TLP)}
Thread-level parallelism executes independent instruction sequences on
separate cores or hardware threads.
\end{definitionbox}
\begin{definitionbox}{Instruction-level parallelism (ILP)}
Instruction-level parallelism exists when one instruction stream contains
independent instructions that a processor can overlap.
\end{definitionbox}

These are different descriptions. A processor with four cores, each running a
different loop iteration, has TLP because four independent instruction
streams execute on separate cores. If each core also contains a five-stage
pipeline, it has ILP because each core overlaps instructions within its own
stream. It is not automatically SIMD: the cores may execute different
instructions or different control paths. A good classification names the
level and the evidence rather than counting only the number of cores.

The rest of this chapter focuses on ILP inside one instruction stream. That is
where hazards become the price of overlap: an instruction may be independent
enough to share the machine with another instruction, or it may need a
resource or value that is not ready yet.

\section{The Five-Stage Pipeline Model}
\label{sec:pipeline}

\subsection{Why Overlap Helps}

Suppose every instruction needs five stages: fetch, decode, execute, memory,
and writeback. If one instruction occupies the whole path before the next
starts, four stage resources sit idle for most of the time. A \key{pipeline}
places pipeline registers between stages and overlaps different instructions
in different stages. After the pipeline fills, one completed instruction can
emerge each cycle even though each individual instruction still requires all
five stages. The promise is throughput, not a magical reduction of every
instruction's latency. This is the standard pipeline motivation in Hamacher,
Ch.~10 \cite{Hamacher2002_computer_organization}, and Stallings 11e, S16.4,
PDF pp.~667--685 \cite{Stallings2015_computer_organization}.

\begin{definitionbox}{Five-stage pipeline}
The teaching model divides an instruction path into IF (instruction fetch),
ID (decode and register read), EX (ALU operation or branch work), MEM (data
memory access), and WB (register write). The same instruction visits all five
stages, while different instructions occupy different stages simultaneously.
\end{definitionbox}

Each stage advances on a clock edge. Pipeline registers hold the intermediate
results and control information needed by the next stage. The slowest stage,
plus the pipeline-register overhead, constrains the clock period. Real
processors may split, merge, or repeat these functions; the five-stage model
is a deliberately useful abstraction for reasoning about timing and hazards.

\subsection{Filling the Pipeline}

\begin{example}[Four independent instructions]
\label{ex:fill}
Assume four independent instructions and one cycle per stage. With no hazards,
the schedule is:
\begin{center}
\begin{tabular}{@{}c|cccccccc@{}}
\toprule
 & 1 & 2 & 3 & 4 & 5 & 6 & 7 & 8 \\
\midrule
$I_1$ & IF & ID & EX & MEM & WB & & & \\
$I_2$ & & IF & ID & EX & MEM & WB & & \\
$I_3$ & & & IF & ID & EX & MEM & WB & \\
$I_4$ & & & & IF & ID & EX & MEM & WB \\
\bottomrule
\end{tabular}
\end{center}

Without overlap, each instruction occupies five stage-cycles, so the total is
$4\times5=20$ stage-cycles. With overlap, the first instruction takes five
cycles to emerge and each of the remaining three takes one additional cycle:
$5+(4-1)=8$ cycles. The finite-stream speedup is therefore
\[
  \frac{20}{8}=2.5.
\]
It is below the ideal five because the stream must first fill the pipeline and
then drain it. The timing table makes both effects visible: cycle 1 contains
only $I_1$'s IF stage, and after $I_4$ enters, the already-started
instructions still need to finish.
\end{example}

\subsection{Latency, Throughput, and Speedup}

Three quantities answer three different questions. \key{Latency} is the time
from an instruction's IF to its WB. \key{Throughput} is how often completed
instructions emerge. \key{Speedup} compares total completion time with a
non-overlapped baseline. Confusing latency with throughput produces the false
claim that a five-stage pipeline makes every instruction take one cycle.

With $k$ balanced stages and $n$ independent instructions, the idealized
pipelined time is
\[
  T_{\mathrm{pipe}}\approx k+(n-1)\text{ cycles},
\]
whereas sequential time is $nk$ cycles. For large $n$, ideal speedup
approaches $k$, provided stages are balanced and hazards, fill, drain, and
pipeline-register overhead are negligible. An individual instruction still
traverses all $k$ stages, so its latency remains about $k$ cycles in this
model. The valid claim is that, after fill, independent instructions can
complete at a rate of one per cycle.

Stalls and flushes reduce achieved throughput. A stall appears as a bubble in
the timing chart; a flush appears as work that was fetched but is discarded.
A deeper pipeline can raise clock frequency by shortening individual stages,
but it can also increase branch penalty, register overhead, and recovery cost.

\subsection{Questions That Shape a Pipeline}

Pipeline design is an architectural choice because it determines CPI, area,
power, and verification burden. The designer must ask:
\begin{enumerate}
  \item \textbf{Stage balance:} can work be divided so no stage dominates the
        clock period?
  \item \textbf{Pipeline registers:} what state must cross each boundary, and
        what setup and clock-to-q overhead is added?
  \item \textbf{Resource replication:} can instruction fetch and data access
        happen together, or must they share one memory port?
  \item \textbf{Dependence handling:} where can a value be forwarded, and
        when must the pipeline wait?
  \item \textbf{Control recovery:} how many younger stages must be flushed
        after a branch is resolved?
\end{enumerate}

\section{Hazard Taxonomy}
\label{sec:hazards}

\begin{definitionbox}{Pipeline hazard}
A pipeline hazard is a condition that prevents the next instruction from
executing in its intended cycle. The remedy must match the cause: add or
duplicate a resource, move or forward data, or predict and recover from
control flow.
\end{definitionbox}

There are three principal categories. A \key{structural hazard} occurs when
two stages need the same hardware at once. A \key{data hazard} occurs when an
instruction depends on a value whose production or writeback has not
completed. A \key{control hazard} occurs when the next instruction address is
uncertain, usually because of a branch or exception.

\subsection{Structural Hazards}

In a five-stage pipeline, suppose $I_1$ is in MEM while $I_2$ is in IF. If one
single-ported memory serves both instruction fetches and data accesses, both
stages request the same port in the same cycle. The problem is resource
demand, not a relationship between the values computed by the instructions.

There are three direct responses:
\begin{enumerate}
  \item Stall one instruction, inserting a bubble.
  \item Split instruction and data memories or caches so the accesses proceed
        in parallel.
  \item Use multiported or time-multiplexed hardware, paying an area, energy,
        or timing cost.
\end{enumerate}

\subsection{Data Hazards and Dependence Names}

The running stream begins:
\begin{center}
\texttt{LW R1,0(R2)}\\
\texttt{ADD R3,R1,R4}\\
\texttt{SUB R5,R3,R6}
\end{center}
The ADD needs the value produced by the LW, so the pair has a read-after-write
(RAW) dependence. The SUB needs the value produced by the ADD, so a second RAW
dependence chains the stream. These are true dependences: the consumer really
needs the producer's value.

Write-after-read (WAR) and write-after-write (WAW) hazards can arise when
instructions execute or complete out of order. In a simple in-order
five-stage pipeline, RAW is the central hazard because the normal stage order
already limits the false name conflicts.

\begin{example}[A load-use stall]
\label{ex:load-use}
Assume a load's data becomes available after MEM, while the dependent ALU
instruction needs it at the start of EX. Even with forwarding, the value is
available one stage too late for the immediately following ADD. The schedule
therefore contains one bubble:
\begin{center}
\begin{tabular}{@{}c|cccccccc@{}}
\toprule
 & 1 & 2 & 3 & 4 & 5 & 6 & 7 & 8 \\
\midrule
\texttt{LW}  & IF & ID & EX & MEM & WB & & & \\
\texttt{ADD} & & IF & ID & \textbf{stall} & EX & MEM & WB & \\
\texttt{SUB} & & & IF & \textbf{stall} & ID & EX & MEM & WB \\
\bottomrule
\end{tabular}
\end{center}
The hazard unit freezes the PC and the IF/ID pipeline register for one cycle,
while inserting a bubble into EX by clearing the control fields entering that
stage. After the wait, the load result can be forwarded from the MEM/WB
boundary to the dependent ADD's EX input. The stall is not a failure of
forwarding; it is the one cycle needed because memory data arrives too late.
\end{example}

\subsection{Control Hazards}

A branch changes the instruction stream. The fetch unit may fetch a fall-through
instruction before the branch condition and target are known. If the branch is
taken, those instructions are wrong-path work and must be prevented from
changing architectural state. The penalty is the number of younger stages
that must be flushed or redirected when the branch resolves.

The designer can delay the decision, which is simple but inserts bubbles; can
predict and recover, which keeps fetching but flushes on a misprediction; or
can move branch resolution earlier, which reduces the penalty but lengthens an
earlier stage and may add hardware to the critical path.

\begin{example}[Diagnosing three hazards]
\label{ex:taxonomy}
\begin{enumerate}
  \item If one memory port is requested by IF and MEM in the same cycle, the
        hazard is structural. Duplicate, split, or schedule the resource.
  \item If an ADD reads a register before a preceding load produces it, the
        hazard is data/RAW. Forward when possible and otherwise stall.
  \item If a conditional branch has not resolved its target, the hazard is
        control. Predict, delay, or flush wrong-path work.
\end{enumerate}
In each case, a complete answer gives both the category and the mechanism that
pays for avoiding it.
\end{example}

\section{Resolving Hazards and Spending the Throughput Budget}
\label{sec:resolution}

\subsection{Stalls and Bubbles}

When the next instruction cannot safely advance, the hazard unit holds the
blocked instructions and inserts a no-op bubble into the following stage. The
program's meaning is preserved because the dependent operation waits for the
required resource or value. A stall can freeze the PC, freeze an inter-stage
register, and clear the control fields entering the next stage.

The cost is extra cycles. If a fraction $s$ of instructions incurs an average
of $b$ bubbles, then, before other effects are included,
\[
  \mathrm{CPI}=1+s\times b.
\]
Thus a hazard remedy should be evaluated by the useful work it preserves and
the cycles it adds, not merely by whether it keeps a stage occupied.

\subsection{Forwarding and Bypassing}

Suppose an ALU result is already at the EX/MEM boundary. The next ALU
operation can consume that result directly instead of waiting for WB to write
the register file and ID to read it. Comparators detect whether a source
register matches a destination register in a later stage, and multiplexers
select the newest available value for each ALU input. Forwarding or data
bypassing removes many ALU-to-ALU RAW stalls. It cannot remove a load-use delay
when memory data arrives too late. Forwarding is part of the indexed pipeline
treatment in Stallings 11e, S16.4, and its RISC discussion is cited in the
deck with the same source \cite{Stallings2015_computer_organization}.

\subsection{Scheduling, Renaming, and Readiness Tracking}

The compiler can move an independent instruction into a delay slot or between
a producer and consumer when the ISA permits it. Register renaming gives
independent writes different physical destinations, removing false WAR and
WAW name conflicts. Scoreboarding and reservation stations track operand
readiness and resource availability so independent instructions can proceed.

These techniques expose more ILP, but require metadata, selection logic,
recovery rules, and stronger verification. No scheduling trick can remove a
true RAW dependence: the consumer still needs the producer's value.

\subsection{Prediction and Recovery}

A branch predictor guesses direction and sometimes the target so that fetch
continues without waiting for full branch resolution. A correct prediction
turns a control hazard into useful work. A wrong prediction requires a
\key{flush}: younger instructions are invalidated and fetch restarts from the
correct target. The expected branch cost is approximately
\[
  \Delta\mathrm{CPI}\approx f_{br}\,p_{miss}\,P_{flush},
\]
where $f_{br}$ is branch frequency, $p_{miss}$ is the misprediction rate, and
$P_{flush}$ is the recovery penalty. Deeper pipelines can increase
$P_{flush}$ even when their clock is faster.

\subsection{A Design Comparison}

Every architecture choice spends some part of a throughput budget.
\begin{center}
\begin{tabular}{@{}>{\raggedright\arraybackslash}p{2.8cm}>{\raggedright\arraybackslash}p{3.2cm}>{\raggedright\arraybackslash}p{3.2cm}>{\raggedright\arraybackslash}p{3.0cm}@{}}
\toprule
\textbf{Choice} & \textbf{Potential gain} & \textbf{New cost} &
\textbf{Question to ask} \\
\midrule
Deeper pipeline & Shorter stage delay; higher clock & More registers, bubbles, branch penalty & Are hazards predictable enough? \\
Duplicate resources & Fewer structural conflicts & Area, power, coherence/consistency & Is the conflict frequent? \\
Forwarding network & Fewer RAW stalls & Comparators, multiplexers, timing paths & Can the value arrive in time? \\
Branch prediction & Higher control-flow throughput & Predictor state and flush recovery & What is the miss penalty? \\
Out-of-order issue & More ILP exposed & Rename, queues, precise recovery & Is added complexity worth it? \\
\bottomrule
\end{tabular}
\end{center}

The design principle is to maximize useful completed work per unit time, not
merely the clock rate or the number of stages.

\begin{example}[A complete diagnosis of the running stream]
\label{ex:running-stream}
Consider again:
\begin{center}
\texttt{LW R1,0(R2)}; \texttt{ADD R3,R1,R4}; \texttt{SUB R5,R3,R6};
\texttt{BEQ R5,R0,L}; \texttt{OR R7,R8,R9}.
\end{center}
The LW to ADD edge is a true RAW dependence. In the five-stage timing model,
it requires forwarding plus one load-use stall. The ADD to SUB edge is another
true RAW dependence, but ALU forwarding can avoid an additional stall. The SUB
to BEQ edge means that the branch depends on a just-produced comparison value,
so forwarding or delayed branch resolution matters. Finally, the BEQ to OR
edge is a control hazard: fetch may need prediction and a flush if the branch
is taken or mispredicted.

A timing diagram is therefore a proof of the schedule. For each dependence,
mark the producer stage, consumer stage, available bypass path, and every
inserted bubble. This prevents the common mistake of saying ``forwarding
solves hazards'' without checking when the value actually becomes available.
\end{example}

\subsection{Choosing the Cheapest Correct Fix}

If an ALU result is ready at the end of EX, its consumer enters EX in the next
cycle, and an ALU bypass path already exists, the hazard unit should forward
the result; no bubble is needed. If the producer is a load, the hazard unit
must stall until the data is available and then forward it. The correct
resolution depends on when the value becomes available, not only on the fact
that a dependence exists.

\section{Decision Procedure and Week 12 Bridge}
\label{sec:synthesis}

For an unfamiliar parallel or pipeline problem, use the following procedure.
\begin{enumerate}
  \item Identify the streams: count the instruction streams and data streams.
  \item For a pipeline conflict, name the blocked stage and the requested
        resource or operand.
  \item Classify the hazard as structural, data, or control.
  \item Ask whether duplication, forwarding, scheduling, prediction, or a
        stall resolves it without changing program meaning.
  \item Count the bubbles or flushes and update the timing diagram.
  \item State the trade-off in throughput, latency, area, power, complexity,
        and recovery cost.
\end{enumerate}

The central discipline is to connect a label to evidence. ``Many ALUs'' does
not establish SIMD; a dependence does not by itself determine whether a stall
is needed; and a faster clock does not by itself establish higher throughput.
The organization, timing, and available resources must all be named.

Next week turns these mechanisms into worked instruction-pipelining problems
and compares RISC and CISC choices that make hazards easier or harder to
manage.

\subsection*{Exercises}
\begin{enumerate}
  \item Classify each workload as SISD, SIMD, MISD, or MIMD and justify the
        classification by naming its instruction and data streams: (a) one
        scalar core sums an array; (b) one vector instruction adds eight array
        elements; (c) four cores run independent matrix tiles.
  \item For six independent instructions in a balanced five-stage pipeline,
        compute the sequential time, ideal pipelined time, and finite-stream
        speedup. State the limiting large-stream speedup.
  \item Diagnose each hazard and give its cheapest correct response: (a) IF
        and MEM request one single-ported memory; (b) an ALU result is ready at
        the end of EX and its consumer enters EX next cycle; (c) a load's data
        is available only after MEM.
  \item For the running stream, identify every true dependence and the
        control hazard. State which dependence needs one load-use bubble and
        which can be handled by ALU forwarding.
  \item A workload has branch frequency $f_{br}=0.20$, misprediction rate
        $p_{miss}=0.05$, and flush penalty $P_{flush}=3$ cycles. Compute the
        approximate branch contribution to CPI and explain what happens if a
        deeper pipeline raises the penalty to 5 cycles.
\end{enumerate}

\subsection*{Solutions}
\begingroup\small
\begin{enumerate}
  \item (a) SISD: one instruction stream operates on one scalar data stream.
        (b) SIMD: one instruction stream applies the same operation to many
        data elements. (c) MIMD: the cores run independent instruction streams
        over independent data regions.
  \item Sequential time is $6\times5=30$ stage-cycles. The ideal pipelined
        time is $5+(6-1)=10$ cycles. The finite-stream speedup is $30/10=3$.
        For a large independent stream, the ideal speedup approaches 5, not
        1, because one instruction can complete each cycle after fill.
  \item (a) Structural hazard: split or duplicate the memory if the conflict
        justifies its area and power; otherwise insert a bubble. (b) Data/RAW
        hazard with a value available in time: forward the ALU result, with no
        bubble. (c) Data/RAW load-use hazard: stall until MEM produces the
        value, then forward it.
  \item The LW produces $R1$, which ADD reads: true RAW and one load-use stall.
        ADD produces $R3$, which SUB reads: true RAW handled by ALU
        forwarding. SUB produces the comparison value used by BEQ: another
        true RAW whose handling depends on branch-resolution timing. BEQ
        changes the next instruction address, so the following OR is exposed
        to a control hazard and may be flushed after a wrong prediction.
  \item The branch contribution is
        \[
          \Delta\mathrm{CPI}\approx0.20\times0.05\times3=0.03.
        \]
        If the penalty rises to 5 cycles, it becomes
        $0.20\times0.05\times5=0.05$ CPI. The clock may be faster, but each
        misprediction costs more cycles, so useful throughput depends on both
        effects.
\end{enumerate}
\endgroup

\section*{References and Source Boundaries}

The deck's references are retained without changing their scope. Hamacher,
Vranesic, Zaky, and Manjikian, \emph{Computer Organization and Embedded
Systems}, 6th ed., Chs.~5 and 10, is a chapter-level course anchor; the
current repository index does not page-verify these chapters
\cite{Hamacher2002_computer_organization}. Stallings, \emph{Computer
Organization and Architecture: Designing for Performance}, 11th ed. (the
repository key is retained as Stallings2015), is cited for S16.4,
``Instruction Pipelining,'' PDF pp.~667--685, and S20.1, ``Multiple Processor
Organizations,'' PDF p.~847 \cite{Stallings2015_computer_organization}.
Flynn's taxonomy is not attributed to Stallings here: its treatment in the
indexed 11th edition is brief, so the course anchor remains Hamacher Ch.~5.
Page numbers are used only where the deck itself marks them as checkable.

\begingroup
\raggedright
\bibliographystyle{plain}
\bibliography{../../Bibliography_base}
\endgroup

\end{document}